% ====================================================================
% SECTION 09: MECHANICAL-FLUIDIC DECOUPLING
% ====================================================================
\section{The Retinal Paradox: Mechanical-Fluidic Decoupling}
\label{sec:retinal}

The human retina provides a unique test case for the incommensurability
principle. Unlike three-dimensional vascular beds (coronary, pulmonary), retinal
vessels are constrained to a quasi-two-dimensional manifold (the retinal
surface). According to the dual-threshold framework derived from the Kinematic
Matching Criterion (Theorem~\ref{thm:kinematic_matching},
Section~\ref{sec:womersley_minimax}), this dimensional constraint produces a
profound qualitative shift:

\subsection{Predicted Dimensional Scaling}

The fluid and wave thresholds scale differently with dimension $d$:
\begin{itemize}
    \item
\textbf{Fluid threshold (local admittance $Y_L$):}
$\mathrm{Wo}_c^{\mathrm{fluid}}(d) = \sqrt{6/(d-1)}$

For $d=3$: $\mathrm{Wo}_c^{\mathrm{fluid}} = \sqrt{3} \approx 1.732$

For $d=2$: $\mathrm{Wo}_c^{\mathrm{fluid}} = \sqrt{6} \approx \VarWoCTwo$

    \item
\textbf{Wave threshold (characteristic admittance $Y_c \propto \sqrt{Y_L}$):}
$\mathrm{Wo}_c^{\mathrm{wave}}(d) = \sqrt{3d(d-2)/(d-1)}$

For $d=3$: $\mathrm{Wo}_c^{\mathrm{wave}} = 3/\sqrt{2} \approx 2.121$

For $d=2$: $\mathrm{Wo}_c^{\mathrm{wave}} = 0$ (collapses identically)
\end{itemize}

The collapse of $\mathrm{Wo}_c^{\mathrm{wave}}$ to zero in planar topology
implies that \emph{any} finite Womersley number in 2D networks places the wave
permanently above its kinematic matching threshold ($\mathcal{Q}^{-1}_{Y_c} \ge
1$ for all $\mathrm{Wo} > 0$). This geometric constraint forces retinal vessels
into a permanent wave-dominated regime.

Because the human retinal vasculature is embedded in a systemic tree operating
well above the 2D transition mass ($M \gg M^*_{d=2}$), the dimensional
constraint shifts the minimax diameter attractor toward the area-preserving
limit. The minimax prediction is:
\begin{itemize}
    \item
\textbf{Diameters:} $\alpha_{\text{retina}} \approx 2.0$ (area-preserving limit)
    \item
\textbf{Angles:} $\varphi_{\text{retina}} \approx 90^\circ$ (planar constraint:
$\tan^2(\varphi/2) = 1$)
\end{itemize}

For comparison, Murray's law in 3D predicts $\alpha_{\mathrm{3D}} \approx 3.0$
(viscous) and $\varphi_{\mathrm{3D}} \approx \VarAngleTetrahedralCalcDeg^\circ$
(derived from the energy-minimization condition $\cos(\varphi/2) = 2^{-1/3}$ for
symmetric bifurcations).

\subsection{Observational Evidence}

\citet{luo2017} measured 342 bifurcations in healthy human retinal vessels:
\begin{itemize}
    \item
\textbf{Diameters:} $\alpha_{\text{obs}} \approx 2.0$--$2.7$ (range across
vessel types; consistent with area-preserving limit)
    \item
\textbf{Angles:} $\varphi_{\text{obs}} = \VarAngleRetinalObs^\circ \pm
\VarAngleRetinalErr^\circ$ (SE), with standard deviation $\sigma =
\VarAngleRetinalSD^\circ$
\end{itemize}

The diameter scaling ($\alpha_{\text{obs}} \approx 2.0$--$2.7$) is consistent
with the area-preserving attractor predicted by the 2D dimensional constraint.
The observed mean angle $\VarAngleRetinalObs^\circ$ lies between the 3D Murray
prediction ($\VarAngleRetinalCalcDeg^\circ$) and the 2D planar limit
($90^\circ$), offset from the latter by only $\VarAngleRetinalOffset /
\VarAngleRetinalSD \approx 0.20\sigma$; however, given the large biological
variance ($\sigma = \VarAngleRetinalSD^\circ$), the angular data do not
discriminate between the two predictions. The high variance likely reflects
stochastic angiogenesis in vivo, where local biochemical gradients (VEGF, shear
stress) obscure the geometric signal at individual junctions.

\subsection{Effective Viscous Limit and the Fåhræus-Lindqvist Effect}

An additional subtlety arises from the reduction of apparent blood viscosity in
small vessels (Fåhræus-Lindqvist effect). In retinal capillaries ($r <
10\,\mu\text{m}$), the effective viscosity drops by $\sim 30\%$, shifting the
effective transport exponent from $\alpha_t = 3.0$ (Murray) to
$\alpha_{\text{eff}} \approx \VarAlphaFahraeusEff$. This places retinal diameter
scaling between the wave attractor ($\alpha_w = 2.0$) and the corrected viscous
limit ($\alpha_{\text{eff}} = \VarAlphaFahraeusEff$), further confirming the
transition-regime interpretation.

Taken together, the retinal data provide a case study in which the diameter
scaling unambiguously supports the dimensional decoupling predicted by the
incommensurability framework, while the angular evidence remains inconclusive
due to biological variance.
